\section{Compiling \acr{GPU} \acr{LLVM} \acr{IR} to Machine Code}
\label{sec:gpu-compilation}

After generating \acr{LLVM} \acr{IR} with \acr{GPU}-specific intrinsics, you must compile it to the target \acr{GPU}'s native format using \texttt{llc} or vendor-specific tools.

\subsection{NVIDIA: Compiling \acr{NVVM} \acr{IR} to \acr{PTX}}

\textbf{\acr{PTX}} is NVIDIA's intermediate assembly language for \acr{CUDA}. It's then \acr{JIT}-compiled by the \acr{GPU} driver to native \acr{SASS} for your specific \acr{GPU} architecture.

\textbf{Using llc to generate \acr{PTX}:}
\begin{lstlisting}[style=ir]
# Basic PTX generation
llc -march=nvptx64 -mcpu=sm_75 -o kernel.ptx module.ll

# With optimizations (recommended)
llc -O3 -march=nvptx64 -mcpu=sm_75 -o kernel.ptx module.ll
\end{lstlisting}

\textbf{Common NVIDIA compute capabilities (-mcpu):}
\begin{itemize}[noitemsep]
  \item \texttt{sm\_50}: Maxwell (GTX 900 series)
  \item \texttt{sm\_60}: Pascal (GTX 10xx, Tesla P100)
  \item \texttt{sm\_70}: Volta (Tesla V100)
  \item \texttt{sm\_75}: Turing (\acr{RTX} 20xx, T4)
  \item \texttt{sm\_80}: Ampere (A100, \acr{RTX} 30xx)
  \item \texttt{sm\_86}: Ampere (\acr{RTX} 30xx mobile)
  \item \texttt{sm\_89}: \acr{Ada} Lovelace (\acr{RTX} 40xx)
  \item \texttt{sm\_90}: Hopper (H100)
\end{itemize}

\textbf{Loading \acr{PTX} from Node.js:}
\begin{lstlisting}[style=ts]
import { readFileSync } from "fs";

// Assuming you have a CUDA runtime binding (e.g., gpu.js, node-cuda)
const ptxCode = readFileSync("kernel.ptx", "utf-8");
// Load PTX module and launch kernel with CUDA runtime
\end{lstlisting}

\subsection{\acr{AMD}: Compiling \acr{AMDGPU} \acr{IR} to \acr{GCN}/\acr{RDNA}}

\acr{AMD} \acr{GPU}s use \textbf{\acr{GCN}} or \textbf{\acr{RDNA}} instruction sets. The \acr{LLVM} \acr{AMDGPU} backend can generate both object files and assembly.

\textbf{Using llc for \acr{AMD} \acr{GPU}s:}
\begin{lstlisting}[style=ir]
# Generate AMD GCN assembly
llc -march=amdgcn -mcpu=gfx900 -o kernel.s module.ll

# Generate object file (for ROCm linking)
llc -O3 -march=amdgcn -mcpu=gfx1030 -filetype=obj -o kernel.o module.ll
\end{lstlisting}

\textbf{Common \acr{AMD} \acr{GPU} architectures (-mcpu):}
\begin{itemize}[noitemsep]
  \item \texttt{gfx803}: Polaris (RX 400/500 series)
  \item \texttt{gfx900}: Vega (Radeon VII, MI25)
  \item \texttt{gfx906}: Vega 20 (MI50, MI60)
  \item \texttt{gfx908}: \acr{CDNA} (MI100)
  \item \texttt{gfx90a}: CDNA2 (MI200 series)
  \item \texttt{gfx1030}: \acr{RDNA}2 (RX 6000 series)
  \item \texttt{gfx1100}: \acr{RDNA}3 (RX 7000 series)
\end{itemize}

\textbf{Integration with \acr{ROCm}:}
\begin{lstlisting}[style=ir]
# Link with ROCm device libraries
clang -target amdgcn-amd-amdhsa -mcpu=gfx906 \
      -o kernel.co module.ll

# Create HSA code object
clang-offload-bundler --type=o \
      --targets=hipv4-amdgcn-amd-amdhsa--gfx906 \
      --inputs=kernel.co --outputs=kernel.hsaco
\end{lstlisting}

\subsection{Intel/\acr{OpenCL}: Compiling to \acr{SPIR-V} Binary}

\textbf{\acr{SPIR-V}} is Khronos's standard intermediate representation for \acr{GPU} compute (\acr{OpenCL}, \acr{Vulkan}).

\textbf{Using llvm-spirv translator:}
\begin{lstlisting}[style=ir]
# First, compile to LLVM bitcode
llc -march=spir64 -filetype=obj -o module.bc module.ll

# Then translate to SPIR-V binary
llvm-spirv module.bc -o kernel.spv
\end{lstlisting}

\textbf{Alternative: Direct \acr{SPIR-V} generation with clang:}
\begin{lstlisting}[style=ir]
clang -cc1 -triple spir64-unknown-unknown \
      -emit-llvm-bc -o module.bc module.ll

llvm-spirv module.bc -o kernel.spv
\end{lstlisting}

\textbf{Loading \acr{SPIR-V} in \acr{OpenCL}:}
\begin{lstlisting}[style=ts]
import { readFileSync } from "fs";

// Using an OpenCL binding for Node.js
const spirvBinary = readFileSync("kernel.spv");
// Create OpenCL program from SPIR-V binary
// const program = clCreateProgramWithIL(context, spirvBinary, ...);
\end{lstlisting}

\subsection{\acr{GPU} Compilation Workflow Summary}

\begin{longtable}{|p{3cm}|p{4cm}|p{6cm}|}
\hline
\textbf{\acr{GPU} Vendor} & \textbf{Output Format} & \textbf{Compilation Command} \\
\hline
NVIDIA & \acr{PTX} (text) & \texttt{llc -march=nvptx64 -mcpu=sm\_XX -o out.ptx} \\
\hline
\acr{AMD} & \acr{GCN} assembly or object & \texttt{llc -march=amdgcn -mcpu=gfxXXXX -o out.s} \\
\hline
Intel/\acr{OpenCL} & \acr{SPIR-V} binary & \texttt{llc -march=spir64 | llvm-spirv -o out.spv} \\
\hline
\end{longtable}

\textbf{Key differences:}
\begin{itemize}[noitemsep]
  \item \textbf{NVIDIA \acr{PTX}} is human-readable text (like assembly), \acr{JIT}-compiled by the driver at runtime
  \item \textbf{\acr{AMD} \acr{GCN}/\acr{RDNA}} can be assembly (.s) or binary (.o), linked with \acr{ROCm} runtime
  \item \textbf{\acr{SPIR-V}} is a binary format, consumed by \acr{OpenCL}/\acr{Vulkan} drivers
\end{itemize}

Always use \texttt{-O3} optimization for \acr{GPU} code—the parallel nature of \acr{GPU} execution benefits heavily from loop unrolling, vectorization, and memory coalescing optimizations.

